%% paper3.tex — Development System: архитектура LLM-системы на Common Lisp
%% Компиляция: lualatex paper3.tex
\documentclass[11pt,a4paper]{article}

\usepackage{fontspec}
\setmainfont{DejaVu Serif}
\setmonofont[Scale=0.88]{DejaVu Sans Mono}

\usepackage{geometry}
\geometry{a4paper, top=25mm, bottom=25mm, left=28mm, right=28mm}

\usepackage{polyglossia}
\setmainlanguage{russian}
\setotherlanguage{english}

\usepackage{microtype}
\usepackage{parskip}
\usepackage{xcolor}
\usepackage{booktabs}
\usepackage{enumitem}
\usepackage{titlesec}
\usepackage{fancyhdr}
\usepackage{hyperref}
\usepackage{float}

% ─── листинги ────────────────────────────────────────────────────────────────
\usepackage{listings}

\definecolor{bg}{HTML}{F7F7F5}
\definecolor{kw}{HTML}{0000CC}
\definecolor{cm}{HTML}{3C763D}
\definecolor{st}{HTML}{A31515}
\definecolor{ln}{HTML}{AAAAAA}

\lstdefinelanguage{CL}{
  morekeywords={defun,defpackage,defvar,defparameter,defmacro,
                let,let*,labels,lambda,cond,when,unless,if,progn,
                handler-case,ignore-errors,restart-case,
                loop,dolist,mapc,mapcar,reduce,remove-if,
                format,load,require,in-package,use-package,export,
                setf,push,values,multiple-value-bind,and,or,not,
                append,list,cons,first,rest,car,cdr,null,
                with-open-file,string=,string-upcase,find-symbol,
                find-package,fboundp,probe-file,ensure-directories-exist,
                merge-pathnames,pathname-name,directory,read,prin1,
                funcall,apply,error,warn,sb-ext:exit,uiop:getenv,
                ql:quickload,asdf:defsystem,asdf:load-system,
                asdf:system-source-directory},
  sensitive=false,
  morecomment=[l]{;},
  morecomment=[s]{\#|}{|\#},
  morestring=[b]",
}

\lstset{
  language=CL,
  backgroundcolor=\color{bg},
  basicstyle=\small\ttfamily,
  keywordstyle=\color{kw}\bfseries,
  commentstyle=\color{cm}\itshape,
  stringstyle=\color{st},
  frame=single, framesep=5pt, rulecolor=\color{gray!25},
  showstringspaces=false,
  breaklines=true, breakatwhitespace=true, tabsize=2,
  numbers=left, numberstyle=\tiny\color{ln}, numbersep=7pt,
  xleftmargin=20pt, xrightmargin=4pt,
  captionpos=b,
}

% ─── оформление ──────────────────────────────────────────────────────────────
\hypersetup{colorlinks=true, linkcolor=black, urlcolor=blue!60!black}

\titleformat{\section}{\large\bfseries}{\thesection.}{0.5em}{}
\titleformat{\subsection}{\normalsize\bfseries}{\thesubsection.}{0.5em}{}

\pagestyle{fancy}
\fancyhf{}
\fancyhead[L]{\small\textit{Development System — архитектура на Common Lisp}}
\fancyhead[R]{\small\thepage}
\renewcommand{\headrulewidth}{0.4pt}

% ─── документ ────────────────────────────────────────────────────────────────
\begin{document}

\begin{titlepage}
\centering\vspace*{2.5cm}
{\LARGE\bfseries Development System\\[0.4em]
\large Архитектура многопотоковой LLM-системы на Common Lisp}

\vspace{1.5cm}
{\large Артемий Уродовских}\\[0.2em]
{\normalsize Март 2026}

\vspace{2cm}
\begin{abstract}
\noindent
\textbf{Development System} — двухкомпонентная система оркестрации вычислительных потоков с языковой моделью в роли исполнителя. \textit{Мастер} — Telegram-бот на Common Lisp, принимающий команды и управляющий потоками. \textit{Подмастерье} — персистентный SBCL-образ, в который потоки загружаются динамически без перезапуска.

Поток — автономный модуль CL, определяющий единственную функцию \texttt{выполнить/1}. Поток-порождатель с помощью LLM генерирует новые потоки, компилирует и загружает их в тот же образ, замыкая метациркулярную петлю.

Система полностью написана на Common Lisp (SBCL). Около 1100 строк кода. Мигрирована с Racket в марте 2026 года.
\end{abstract}

\vspace{1cm}
\noindent\textbf{Ключевые слова:}
Common Lisp, SBCL, LLM, Telegram-бот, персистентный образ, метациркулярность, явное состояние
\end{titlepage}

\tableofcontents
\newpage

% ═══════════════════════════════════════════════════════════════════════════════
\section{Обзор архитектуры}

Система состоит из двух SBCL-процессов и каталога потоков:

\begin{figure}[H]\centering\small\begin{verbatim}
Telegram
   │
   ▼
МАСТЕР (мастер/cl/)          SBCL-процесс #1
   poll-loop · роутинг · душа · история диалогов
   │ load .lisp (один раз, ленивая загрузка)
   ▼
мастер/потоки/*.lisp         вычислительные потоки
   │ load (при первом вызове)
   ▼
ПОДМАСТЕРЬЕ (подмастерье/)   SBCL-процесс #2
   персистентный образ · загружает потоки · вызывает напрямую
\end{verbatim}
\caption{Архитектура системы}\end{figure}

\subsection{Контракт потока}

Каждый поток — CL-файл с одним экспортом:

\begin{lstlisting}[caption={Минимальный поток}]
(defpackage :поток-эхо (:use :cl) (:export #:выполнить))
(in-package :поток-эхо)

(defun выполнить (задача)
  (format nil "Эхо: ~a" задача))
\end{lstlisting}

Один пакет \texttt{:поток-ИМЯ}, одна функция \texttt{выполнить/1}. Этот контракт достаточен для автоматизированного запуска и генерации потоков через LLM.

\subsection{Файловая структура}

\begin{table}[H]\centering\small
\begin{tabular}{ll}
\toprule
\textbf{Файл} & \textbf{Роль} \\
\midrule
\texttt{мастер/cl/мастер.asd}    & ASDF-система, порядок загрузки \\
\texttt{мастер/cl/packages.lisp} & Пакет \texttt{:мастер} \\
\texttt{мастер/cl/config.lisp}   & Переменные окружения и пути \\
\texttt{мастер/cl/telegram.lisp} & HTTP-клиент Telegram API \\
\texttt{мастер/cl/llm.lisp}      & HTTP-клиент OpenRouter \\
\texttt{мастер/cl/dusha.lisp}    & Личность бота (plist S-expr) \\
\texttt{мастер/cl/история.lisp}  & Персистентная история диалогов \\
\texttt{мастер/cl/потоки.lisp}   & Загрузка и запуск потоков \\
\texttt{мастер/cl/комбо.lisp}    & Комбинирование потоков \\
\texttt{мастер/cl/команды.lisp}  & Роутинг команд \\
\texttt{мастер/cl/main.lisp}     & Poll-loop, точка входа \\
\texttt{подмастерье/main.lisp}   & Персистентный исполнитель \\
\texttt{мастер/потоки/*.lisp}    & Вычислительные потоки \\
\bottomrule
\end{tabular}
\end{table}

% ═══════════════════════════════════════════════════════════════════════════════
\section{Мастер}

\subsection{ASDF-система}

\begin{lstlisting}[caption={мастер/cl/мастер.asd}]
(asdf:defsystem #:мастер
  :description "Telegram bot / orchestrator"
  :depends-on  (#:dexador #:cl-json)
  :serial t
  :components ((:file "packages")  (:file "config")
               (:file "telegram")  (:file "llm")
               (:file "dusha")     (:file "история")
               (:file "потоки")    (:file "комбо")
               (:file "команды")   (:file "main")))
\end{lstlisting}

\texttt{:serial t} обеспечивает правильный порядок компиляции: каждый файл видит символы предыдущих.

\subsection{Конфигурация}

Пути вычисляются относительно ASDF-системы, а не рабочей директории:

\begin{lstlisting}[caption={мастер/cl/config.lisp}]
(in-package #:мастер)

(defparameter *токен*   (uiop:getenv "MASTER_BOT_TOKEN"))
(defparameter *admin-id* (uiop:getenv "ADMIN_CHAT_ID"))

;; Корень = мастер/ (родитель мастер/cl/ где лежит .asd)
(defparameter *корень*
  (uiop:pathname-parent-directory-pathname
   (asdf:system-source-directory :мастер)))

(defparameter *путь-души*        (merge-pathnames #P"душа.scm"  *корень*))
(defparameter *каталог-потоков*  (merge-pathnames #P"потоки/"   *корень*))
(defparameter *каталог-истории*  (merge-pathnames #P"история/"  *корень*))
\end{lstlisting}

\subsection{Telegram HTTP-клиент}

\begin{lstlisting}[caption={мастер/cl/telegram.lisp}]
(defparameter *telegram-base* "https://api.telegram.org/bot")

(defun api-url (method)
  (concatenate 'string *telegram-base* *токен* "/" method))

(defun api-post (method payload)
  "POST JSON к Telegram API -> alist или nil."
  (handler-case
      (cl-json:decode-json-from-string
       (dexador:post (api-url method)
                     :content (cl-json:encode-json-to-string payload)
                     :headers '(("Content-Type" . "application/json"))))
    (error (e)
      (format *error-output* "[tg] ~A: ~A~%" method e)
      nil)))

(defun get-updates (&key (timeout 30) offset)
  (let ((p (append `((:timeout . ,timeout))
                   (when offset `((:offset . ,offset))))))
    (cdr (assoc :result (api-post "getUpdates" p)))))

(defun send-message (chat-id text &key reply-markup)
  (api-post "sendMessage"
            (append `((:chat--id . ,chat-id) (:text . ,text))
                    (when reply-markup `((:reply--markup . ,reply-markup))))))
\end{lstlisting}

Ключевое решение: \texttt{api-post} никогда не бросает условие — любая ошибка (сеть, 4xx, 5xx) логируется и возвращается \texttt{nil}. Poll-loop остаётся устойчивым к сетевым сбоям.

Замечание по \texttt{cl-json}: библиотека конвертирует \texttt{snake\_case} JSON-ключи в символы CL с двойным дефисом (\texttt{update\_id} → \texttt{:update--id}, \texttt{callback\_query} → \texttt{:callback--query}).

\subsection{LLM-клиент}

\begin{lstlisting}[caption={мастер/cl/llm.lisp}]
(defparameter *llm-url*   "https://openrouter.ai/api/v1/chat/completions")
(defparameter *llm-model* "openai/gpt-4.1")
(defparameter *llm-key*   (uiop:getenv "OPENROUTER_API_KEY"))

(defun %auth-headers ()
  `(("Authorization" . ,(format nil "Bearer ~a" *llm-key*))
    ("Content-Type"  . "application/json")))

(defun %build-messages (prompt &key system messages)
  (append (when system   `((("role" . "system")  ("content" . ,system))))
          (or messages `((("role" . "user")    ("content" . ,prompt))))))

(defun llm-complete (prompt &key (model *llm-model*) system messages)
  "Запрос к LLM -> строка-ответ. Никогда не бросает."
  (handler-case
      (let* ((body  (cl-json:encode-json-to-string
                     `(("model" . ,model)
                       ("messages" . ,(%build-messages
                                       prompt :system system
                                       :messages messages)))))
             (resp  (dexador:post *llm-url* :content body
                                 :headers (%auth-headers)))
             (json  (cl-json:decode-json-from-string resp))
             (msg   (cdr (assoc :message
                                (first (cdr (assoc :choices json)))))))
        (or (cdr (assoc :content msg)) "[нет ответа]"))
    (error (e)
      (format *error-output* "[llm] ~A~%" e)
      "[ошибка LLM]")))
\end{lstlisting}

\subsection{Управление потоками}

Имя пакета формируется через \texttt{\%pkg-name}: поток ``эхо'' живёт в пакете \texttt{:ПОТОК-ЭХО}. Загрузка идемпотентна — проверяется через \texttt{find-package}:

\begin{lstlisting}[caption={мастер/cl/потоки.lisp}]
(defparameter *активные-потоки* (make-hash-table :test #'equal))

(defun список-потоков ()
  (mapcar #'pathname-name
          (directory (merge-pathnames "*.lisp" *каталог-потоков*))))

(defun загрузить-все-потоки ()
  (mapc #'загрузить-поток
        (directory (merge-pathnames "*.lisp" *каталог-потоков*))))

(defun %pkg-name (имя)
  (string-upcase (concatenate 'string "поток-" имя)))

(defun %найти-выполнить (имя)
  (let* ((pkg (find-package (%pkg-name имя)))
         (sym (when pkg (find-symbol "ВЫПОЛНИТЬ" pkg))))
    (when (and sym (fboundp sym)) (symbol-function sym))))

(defun запустить-поток (имя задача)
  "Загружает (один раз) и вызывает поток ИМЯ с ЗАДАЧЕЙ."
  (let ((путь (merge-pathnames (format nil "~A.lisp" имя)
                               *каталог-потоков*)))
    (when (probe-file путь)
      (unless (find-package (%pkg-name имя))
        (загрузить-поток путь))
      (let ((fn (%найти-выполнить имя)))
        (when fn
          (setf (gethash имя *активные-потоки*) t)
          (funcall fn задача))))))
\end{lstlisting}

\subsection{Роутинг команд}

\begin{lstlisting}[caption={мастер/cl/команды.lisp — парсинг и роутинг}]
(defun %split-words (str)
  (remove-if #'zerop
             (uiop:split-string str :separator '(#\Space))
             :key #'length))

(defun %cmd-args (text)
  (let ((w (%split-words text)))
    (values (if w (string-downcase (first w)) "") (rest w))))

(defparameter *кнопки*
  '(("🔧 Породить поток" . "/породить")
    ("📋 Потоки"          . "/потоки")
    ("📊 Состояние"       . "/состояние")
    ("💬 Диалог"          . "/диалог")))

(defun обработать-команду (chat-id text)
  (let* ((t0 (string-trim '(#\Space #\Newline) text))
         (t1 (or (cdr (assoc t0 *кнопки* :test #'string=)) t0)))
    (multiple-value-bind (cmd args) (%cmd-args t1)
      (let ((fn (cdr (assoc cmd *команды* :test #'string=))))
        (if fn
            (handler-case (funcall fn chat-id args)
              (error (e) (format nil "Ошибка: ~A" e)))
            (format nil "Неизвестная команда: ~A" cmd))))))
\end{lstlisting}

Обработчик диалога — единственная функция, которой важна история. После рефакторинга история читается один раз:

\begin{lstlisting}[caption={мастер/cl/команды.lisp — /диалог}]
(defun %обр-диалог (chat-id args)
  (when (null args)
    (return-from %обр-диалог "Использование: /диалог <текст>"))
  (let* ((текст   (format nil "~{~A~^ ~}" args))
         (душа    (ignore-errors (читать-душу *путь-души*)))
         (история (добавить-сообщение chat-id "user" текст))
         (ответ   (llm-complete текст
                                :system (when душа
                                          (душа->системный-промпт душа))
                                :messages история)))
    (добавить-сообщение chat-id "assistant" ответ)
    ответ))
\end{lstlisting}

\subsection{Poll-loop}

\begin{lstlisting}[caption={мастер/cl/main.lisp}]
(defvar *polling* t)

(defun start ()
  (format t "[мастер] Загружаем потоки...~%")
  (ensure-directories-exist *каталог-истории*)
  (загрузить-все-потоки)
  (format t "[мастер] Запущен. Ожидаю сообщения...~%")
  (poll-loop 0))

(defun poll-loop (start-offset)
  (loop with offset = start-offset
        while *polling*
        do (let* ((updates (ignore-errors
                             (get-updates :offset offset :timeout 30)))
                  (max-id  offset))
             (dolist (upd (or updates '()))
               (let ((id (cdr (assoc :update--id upd))))
                 (when (and id (> id max-id)) (setf max-id id)))
               (ignore-errors (обработать-update upd)))
             (when (> max-id offset)
               (setf offset (1+ max-id))))))

(defun %chat-id-of (msg)
  (cdr (assoc :id (cdr (assoc :chat msg)))))

(defun обработать-update (upd)
  (cond
    ((assoc :message upd)
     (let* ((msg  (cdr (assoc :message upd)))
            (text (cdr (assoc :text msg))))
       (when (stringp text)
         (send-message (%chat-id-of msg)
                       (обработать-команду (%chat-id-of msg) text)))))
    ((assoc :callback--query upd)
     (let* ((cb   (cdr (assoc :callback--query upd)))
            (cid  (%chat-id-of (cdr (assoc :message cb))))
            (data (cdr (assoc :data cb))))
       (answer-callback-query (cdr (assoc :id cb)))
       (send-message cid (обработать-команду cid data))))))
\end{lstlisting}

Решение против переполнения стека: \texttt{loop} вместо хвостовой рекурсии. SBCL не гарантирует TCO, а \texttt{loop} в CL — полноценная итерация без расхода стека.

% ═══════════════════════════════════════════════════════════════════════════════
\section{Подмастерье}

\subsection{Явное состояние}

Весь Подмастерье написан в стиле SICP: состояние не хранится в глобальных переменных, а передаётся как параметр рекурсии. Реестр загруженных потоков — это \texttt{alist}, который «живёт» в параметрах функций:

\begin{lstlisting}[caption={подмастерье/main.lisp — рекурсивный цикл с явным состоянием}]
(defun опрос (сдвиг реестр)
  (handler-case
      (let* ((тело       (получить-обновления сдвиг))
             (обновления (or (разобрать-обновления тело) '()))
             (сообщения  (извлечь-сообщения обновления)))
        (multiple-value-bind (новый-рег новый-сдвиг)
            (обр-сообщения сообщения реестр сдвиг)
          (when (> новый-сдвиг сдвиг)
            (сохранить-сдвиг новый-сдвиг))
          (опрос новый-сдвиг новый-рег)))
    (error (e)
      (format t "Ошибка цикла: ~a~%" e)
      (sleep 5)
      (опрос сдвиг реестр))))
\end{lstlisting}

\texttt{реестр} — alist вида \texttt{("имя-пакета" . "путь")}. Он накапливается по мере загрузки потоков и нигде не хранится глобально.

\subsection{Ленивая загрузка в образ}

\begin{lstlisting}[caption={подмастерье/main.lisp — идемпотентная загрузка}]
(defun загрузить-в-образ (путь пакет реестр)
  "Загружает поток если ещё не в реестре. -> новый реестр."
  (if (assoc пакет реестр :test #'string=)
      реестр
      (handler-case
          (progn
            (load путь :verbose nil :print nil)
            (format t "Загружен поток: ~a~%" пакет)
            (cons (cons пакет путь) реестр))
        (error (e)
          (format t "Ошибка загрузки ~a: ~a~%" пакет e)
          реестр))))
\end{lstlisting}

Если поток уже есть в реестре — \texttt{load} не вызывается. Второй вызов потока — это просто вызов \texttt{funcall} на уже загруженную функцию.

\subsection{Прямой вызов функции}

\begin{lstlisting}[caption={подмастерье/main.lisp — вызов без subprocess}]
(defun вызвать-поток (пакет задача)
  "Вызывает ВЫПОЛНИТЬ в загруженном пакете напрямую."
  (handler-case
      (let* ((pkg (find-package (string-upcase пакет)))
             (sym (when pkg (find-symbol "ВЫПОЛНИТЬ" pkg))))
        (if (and sym (fboundp sym))
            (funcall sym задача)
            (format nil "Поток ~a не найден." пакет)))
    (error (e) (format nil "Ошибка: ~a" e))))
\end{lstlisting}

Ключевое преимущество: нет сериализации, нет IPC, нет запуска нового процесса. Поток вызывается как обычная CL-функция.

\subsection{Обработка сообщений}

\begin{lstlisting}[caption={подмастерье/main.lisp — рекурсивная обработка списка}]
(defun обр-сообщения (список реестр макс)
  "Рекурсивно обрабатывает список. -> (values реестр макс)."
  (if (null список)
      (values реестр макс)
      (let* ((с       (car список))
             (ид-обн  (first с))
             (ид-чата (second с))
             (текст   (third с)))
        (format t "< ~a~%" текст)
        (multiple-value-bind (новый-рег)
            (if (админ? ид-чата)
                (multiple-value-bind (ответ р)
                    (обработать текст ид-чата реестр)
                  (format t "> ~a~%" ответ)
                  (послать-сообщение ид-чата ответ)
                  р)
                (progn (format t "! чужой: ~a~%" ид-чата) реестр))
          (обр-сообщения (cdr список) новый-рег
                          (max макс (1+ ид-обн)))))))
\end{lstlisting}

Вместо мутабельного цикла — хвостовая рекурсия. Реестр и сдвиг явно передаются через \texttt{values}.

% ═══════════════════════════════════════════════════════════════════════════════
\section{Потоки}

\subsection{Счётчик (тривиальный поток)}

\begin{lstlisting}[caption={мастер/потоки/счетчик.lisp}]
(defpackage :поток-счетчик (:use :cl) (:export #:выполнить))
(in-package :поток-счетчик)

(defun факториал (n)
  (if (<= n 1) 1 (* n (факториал (- n 1)))))

(defun выполнить (задача)
  (handler-case
      (write-to-string (факториал (parse-integer задача)))
    (error () "Ошибка: введите целое неотрицательное число")))
\end{lstlisting}

\subsection{Уточнение (итеративный LLM-поток)}

Поток задаёт до 5 итераций уточнения ответа. Рекурсия с явным счётчиком шага — классический SICP-паттерн:

\begin{lstlisting}[caption={мастер/потоки/уточнение.lisp}]
(defun %готово-ли (ответ)
  (and (> (length ответ) 20)
       (or (search "Ответ:" ответ)
           (search "Итог:"  ответ)
           (search "ГОТОВО" ответ))))

(defun %уточнить (задача предыдущий шаг)
  (if (> шаг 5)
      предыдущий
      (let* ((вопрос (if (string= предыдущий "")
                         (format nil "Задача: ~a" задача)
                         (format nil "Задача: ~a~%Ответ:~%~a~%Уточни."
                                 задача предыдущий)))
             (ответ  (%запрос (list (%системный)
                                    (%сообщение "user" вопрос)))))
        (if (%готово-ли ответ)
            ответ
            (%уточнить задача ответ (1+ шаг))))))

(defun выполнить (задача)
  (handler-case (%уточнить задача "" 1)
    (error (e) (format nil "Ошибка: ~a" e))))
\end{lstlisting}

\subsection{Комбинирование потоков}

Последовательный режим — pipeline через \texttt{reduce}. Параллельный — \texttt{mapcar} с объединением результатов:

\begin{lstlisting}[caption={мастер/cl/комбо.lisp}]
(defun %join (sep strings)
  (reduce (lambda (a b) (concatenate 'string a sep b))
          strings :initial-value ""))

(defun комбинировать (потоки задача &key параллельно)
  (if параллельно
      (%join "\n---\n"
             (mapcar (lambda (имя)
                       (or (ignore-errors (запустить-поток имя задача)) ""))
                     потоки))
      (reduce (lambda (вход имя)
                (or (ignore-errors (запустить-поток имя вход)) вход))
              потоки :initial-value задача)))
\end{lstlisting}

В последовательном режиме \texttt{reduce} элегантно реализует pipeline: результат каждого потока становится входом следующего. Весь паттерн вмещается в голове.

% ═══════════════════════════════════════════════════════════════════════════════
\section{Поток-порождатель}

Замыкание метациркулярной петли: LLM генерирует CL-код → компилируется → загружается в тот же образ:

\begin{lstlisting}[caption={мастер/потоки/порождатель.lisp — точка входа}]
(defun выполнить (описание)
  (handler-case
      (let* ((код   (извлечь-из-блока (запросить-генерацию описание)))
             (имя   (имя-пакета код))
             (путь  (записать-поток имя код))
             (ok?   (компилировать путь)))
        (when ok? (загрузить-в-образ путь (format nil "поток-~a" имя)))
        (format nil "Поток ~s ~a." имя (if ok? "готов" "не скомпилирован")))
    (error (e)
      (format nil "Ошибка порождения: ~a" e))))
\end{lstlisting}

Компиляция проверяется через третье возвращаемое значение \texttt{compile-file}:

\begin{lstlisting}[caption={мастер/потоки/порождатель.lisp — компиляция}]
(defun компилировать (путь)
  (handler-case
      (multiple-value-bind (_ warn fail)
          (compile-file путь :print nil :verbose nil)
        (declare (ignore _ warn))
        (not fail))
    (error () nil)))
\end{lstlisting}

% ═══════════════════════════════════════════════════════════════════════════════
\section{История диалогов}

История хранится в \texttt{история/<chat-id>.json} — список сообщений в формате OpenAI:

\begin{lstlisting}[caption={мастер/cl/история.lisp}]
(defun %путь-истории (chat-id)
  (merge-pathnames (format nil "~A.json" chat-id) *каталог-истории*))

(defun загрузить-историю (chat-id)
  (let ((файл (%путь-истории chat-id)))
    (when (probe-file файл)
      (handler-case
          (cl-json:decode-json-from-string (uiop:read-file-string файл))
        (error () nil)))))

(defun сохранить-историю (chat-id история)
  (let ((файл (%путь-истории chat-id)))
    (ensure-directories-exist файл)
    (with-open-file (out файл :direction :output :if-exists :supersede)
      (write-string (cl-json:encode-json-to-string история) out))))

(defun добавить-сообщение (chat-id role content)
  (let ((новая (append (or (загрузить-историю chat-id) '())
                       (list `((:role . ,role) (:content . ,content))))))
    (сохранить-историю chat-id новая)
    новая))

(defun очистить-историю (chat-id)
  (let ((файл (%путь-истории chat-id)))
    (when (probe-file файл) (delete-file файл))))
\end{lstlisting}

% ═══════════════════════════════════════════════════════════════════════════════
\section{Ключевые решения}

\subsection{Единый язык}

Переход с Racket на CL устранил семантическое несоответствие: оркестратор и потоки в одном пространстве имён. Общие библиотеки (\texttt{dexador}, \texttt{cl-json}) загружаются один раз. Ошибки в потоках ловятся системой условий CL.

\subsection{Идемпотентность через рефлексию}

\texttt{find-package} как ключ в реестре: если пакет уже существует — поток загружен. Не нужен отдельный реестр-переменная. Образ сам является реестром.

\subsection{Явное состояние в Подмастерье}

Реестр потоков и сдвиг Telegram не хранятся глобально. Они передаются параметром через рекурсию \texttt{опрос/2}. Это исключает гонку состояний и делает поведение предсказуемым: перезапуск всегда начинается с чистого реестра.

\subsection{Функции до 15 строк}

Жёсткое ограничение: функция должна умещаться на экране. Если не умещается — нужно выделить вспомогательную функцию. Результат: весь модуль \texttt{telegram.lisp} — 35 строк. Модуль \texttt{llm.lisp} — 25 строк.

\subsection{Метациркулярность}

Порождатель использует \texttt{compile-file} + \texttt{load} — штатные CL-операции. Система расширяет себя в рантайме без перезапуска. Это не хак, а следствие гомоиконности языка.

% ═══════════════════════════════════════════════════════════════════════════════
\section{Запуск}

\begin{lstlisting}[language=bash,caption={start.sh}]
#!/usr/bin/env bash
exec sbcl --noinform \
  --load ~/quicklisp/setup.lisp \
  --eval "(push #p\"мастер/cl/\" asdf:*central-registry*)" \
  --eval '(asdf:load-system :мастер :verbose nil)' \
  --eval '(мастер:start)'
\end{lstlisting}

\begin{table}[H]\centering\small
\begin{tabular}{lll}
\toprule
\textbf{Переменная} & \textbf{Компонент} & \textbf{Значение} \\
\midrule
\texttt{MASTER\_BOT\_TOKEN}     & Мастер      & Telegram Bot API токен \\
\texttt{APPRENTICE\_BOT\_TOKEN} & Подмастерье & Telegram Bot API токен \\
\texttt{ADMIN\_CHAT\_ID}        & Оба         & Whitelist chat ID \\
\texttt{OPENROUTER\_API\_KEY}   & Оба         & Ключ OpenRouter \\
\bottomrule
\end{tabular}
\caption{Переменные окружения}
\end{table}

% ═══════════════════════════════════════════════════════════════════════════════
\section{Заключение}

Development System реализует простую и умещающуюся в голове архитектуру:

\begin{itemize}[noitemsep]
  \item \textbf{Контракт потока} — один пакет, одна функция. Этого достаточно.
  \item \textbf{Персистентный образ} — потоки загружаются один раз и вызываются как функции.
  \item \textbf{Явное состояние} — реестр и сдвиг передаются параметром, не хранятся глобально.
  \item \textbf{Метациркулярность} — система порождает себя через LLM и \texttt{compile-file}.
  \item \textbf{Единый язык} — весь стек на Common Lisp, без переключения контекста.
\end{itemize}

Общий объём — около 1100 строк. Каждый файл умещается на одном экране.

\end{document}
